\scanpage{4}
% archive sheet number: 3
% lines 1–2
Soit $N$ un autre module de prés. finie, Alors on a diagramme commutatif à lignes exactes

% lines 3–4
\diagram{\begin{tikzcd}[ampersand replacement=\&,column sep=small,row sep=small]
0\arrow[r,dash]\&\mathrm{Hom}(M,N)\arrow[r,dash]\arrow[d,"\alpha"]\&N^q\arrow[r,dash]\arrow[d,dash,"\beta"]\&N^p\arrow[d,"\gamma"]\\
\deleted{\unsure{A}}\,0\arrow[r,dash]\&\mathrm{Hom}(\widetilde M,\widetilde N)\arrow[r,dash]\&\Gamma\widetilde N^q\arrow[r,dash]\&\Gamma\widetilde N^p
\end{tikzcd}}

% lines 5–8
et comme $\beta$ et $\gamma$ sont bijectives, de $=$ $\alpha$, donc le foncteur \inserted{$M\unsure{\leadsto|\sim}\widetilde M$} est \uline{pleinement fidèle}. Prouvons qu'il est essentiellement surjectif : en effet, si un $F$

% lines 9–10
coh. sur $X$, on sait qu'il est engendré par ses sections, i.e. $\exists\,\underline O_X^{\deleted{\unsure{p}}q}\to F\to0$ \deleted{\unsure{surjectif}} \inserted{exact}

% lines 11–13
Soit $R$ le noyau, Alors il existe de $=$ $\underline O_X^p - R - 0$, d'où une suite exacte

% lines 14
\[
 \underline O_X^{p}\mathrel{\overset{u}{-}}\underline O_X^{\deleted{\unsure{p}}q} - F - 0
\]

% lines 15–16
et par suite \deleted{\unsure{si} $u$ est donné par} \deleted{et par suite} si $u$ définit. (passant aux sections)

% lines 17
\[
 A^{\deleted{\unsure{q}}p}\mathrel{\overset{u}{\to}}A^q
\]

% lines 18–19
on prend $M=\operatorname{Coker}u$ et on trouve que $F\simeq\widetilde M$.
